\documentclass[10pt]{article}
\usepackage[letterpaper,hmargin=1.02in,vmargin=0.85in]{geometry}
\linespread{1.04}
\usepackage{amsmath,amssymb,amsthm}
\usepackage[hidelinks]{hyperref}
\newtheorem{theorem}{Theorem}
\newcommand{\E}{\mathbb E}
\newcommand{\F}{\mathbb F}
\newcommand{\Z}{\mathbb Z}
\newcommand{\ee}{\mathrm e}
\newcommand{\abs}[1]{\left|#1\right|}
\newcommand{\norm}[1]{\left\|#1\right\|}
\setlength{\parindent}{1.25em}
\title{Sharp cancellation for bilinear sums on weighted trees}
\author{}
\date{}
\begin{document}
\maketitle
\vspace{-3.7em}

\begin{abstract}
Attach a variable to each vertex of a tree and a bilinear phase to each
edge. How large can the average phase become if one may change the phase
at every vertex by an arbitrary function of its variable? Over
$\Z/p^k\Z$ the exact answer is a maximum-weight path problem, with each
edge weighted by the conductor of its character. The upper estimate uses
iterated Fourier transforms. An integral dual solution gives vertex
phases attaining it. Over odd finite fields we also compute the largest
collision probability after arbitrary one-variable perturbations.
\end{abstract}

\section{An exact norm}
The starting point is the path phase $\sum_i x_ix_{i+1}$ and iterated
Parseval in Alp\"oge's account of his proposed RH argument
\cite[p.~3, footnote 3]{alpoge}. We ask what survives on branching paths,
with singular edge coefficients and worst-case local perturbations.
The proofs here use finite Fourier analysis and bipartite linear
programming.

Let $T=(V,E)$ be a finite nonempty tree, $p$ a prime, $k\ge1$, and
$R=\Z/p^k\Z$. Write $\ee_m(t)=\exp(2\pi i t/m)$. Choose $a_e\in R$
on each edge. Write $p^{w_e}$ for the additive order of $a_e$ in $R$,
with $w_e=0$ for $a_e=0$. Define
\[
 W(T,a)=\max_{\substack{L\subseteq E\\ \deg_L(v)\le2\ (v\in V)}}
                \sum_{e\in L}w_e.
\]
Each feasible $L$ is a disjoint union of paths, with isolated vertices
allowed. All expectations below are uniform over the indicated finite
sets and all input coordinates are independent.

\begin{theorem}\label{thm:ring}
For arbitrary complex vertex functions $f_v$ of modulus one,
\[
 \max_{(f_v)}
 \abs{\E_{x\in R^V}
  \ee_{p^k}\!\left(\sum_{uv\in E}a_{uv}x_ux_v\right)
  \prod_{v\in V}f_v(x_v)}
     =p^{-W(T,a)/2}.
 \tag{1}
\]
The same maximum holds for $|f_v|\le1$. For odd $p$ it is attained with
every $f_v$ taking values in the $p^k$-th roots of unity.
\end{theorem}

For example, on a star over $\Z/27\Z$ with coefficients $1,3,9,9$,
the two heaviest spokes have weights $3$ and $2$. The exact maximum
is $3^{-5/2}$. Any number of additional spokes of weight at most two
gives the same maximum. On a path every edge contributes its full weight.

\section{Upper bound by eliminating paths}
For $a\in R$ of order $d=p^w$, put
\[
 K_a g(x)=\E_{y\in R}\ee_{p^k}(axy)g(y).
\]
On $L^2(R)$ with normalized counting measure, $K_a^*K_a$ is
$d^{-1}$ times the orthogonal projection onto functions constant on
cosets of $dR$. Indeed, averaging the character in $x$ gives
$\mathbf1_{y-z\in dR}$. Thus $\norm{K_a}_{2\to2}=d^{-1/2}$,
including the case $a=0$. A normalized sum on a path is a matrix
coefficient of a product of these operators, alternating with
multiplication by the interior vertex functions. The endpoint functions
have $L^2$ norm at most one and the multipliers have operator norm at
most one. Its modulus is at most $\prod_{e\text{ on path}}p^{-w_e/2}$.
For a singleton the bound is one.

Fix any feasible $L$ and contract each component of $(V,L)$ to a point.
The result is a tree. A leaf component therefore meets the rest through
one edge $uv$, with $u$ inside and $v$ outside. For fixed $x_v$, its
edge phase is just another multiplier at $u$. Averaging the leaf
component gives a function of $x_v$ bounded by the path estimate,
uniformly in $x_v$. Divide out that bound, absorb the resulting function
into $f_v$, and repeat. The last path has the same estimate. Multiplying
the bounds and maximizing the retained weight proves the upper bound
in (1).

\section{Attainment from the dual problem}
Write $t_+=\max(t,0)$. The required combinatorial identity is
\[
 W(T,a)=\min_{s_v\in\{0,\ldots,k\}}
 \left(2\sum_v s_v+\sum_{uv\in E}(w_{uv}-s_u-s_v)_+\right).
 \tag{2}
\]
Indeed, relax the indicators of the retained edges to $0\le z_e\le1$
and impose $\sum_{e\ni v}z_e\le2$. The vertex-edge incidence matrix of
a tree is totally unimodular because changing the signs of rows in one
bipartition class gives a directed incidence matrix. Appending unit
rows, transposing, and changing signs preserve total unimodularity.
Hence both this primal program and its dual have integral optimal
vertices for the integral data here \cite{schrijver}. The dual minimizes
$2\sum_v s_v+\sum_e t_e$ subject to $s_u+s_v+t_e\ge w_e$ and
$s_v,t_e\ge0$. Eliminating $t_e$ gives (2), since lowering any
$s_v>k$ to $k$ preserves feasibility.

Choose a minimizer $s$ and first allow $|f_v|\le1$. Support $f_v$ on
$p^{s_v}R$, whose density is $p^{-s_v}$. Conditional on this event
we may write $x_v=p^{s_v}y_v$ with $y_v$ uniform on $R$. Now edge $uv$
has coefficient $b_e=p^{s_u+s_v}a_e$, of order
$d_e=p^{(w_e-s_u-s_v)_+}$.

For any $b\in R$ of order $d$ use an even, unimodular function $c_b$ on
$\Z/d\Z$ satisfying
\[
 c_b(y-z)=c_b(y)c_b(z)\ee_{p^k}(byz),\qquad
 \abs{\E_y c_b(y)}=d^{-1/2}. \tag{3}
\]
Here is an explicit construction. Set $c_0=1$. Otherwise write
$b/p^k=u/d$ modulo one with $u$ coprime to $d$. For odd $d$ take
$c_b(y)=\ee_d(-u\,2^{-1}y^2)$ with inverse in $\Z/d\Z$. For even $d$
take $c_b(y)=\exp(-\pi i u y^2/d)$, which has period $d$ since $d$ is
even. Expanding the square gives the first identity. For the second,
square the modulus of the average, write the two arguments as $y+h,y$,
and sum first in $y$. Orthogonality leaves only $h=0$ modulo $d$.

On the chosen subgroup define
\[
 f_v(p^{s_v}y_v)=\prod_{e\ni v}c_{b_e}(y_v).
\]
This is well defined because every incident $d_e$ divides $p^{k-s_v}$.
By (3) the whole integrand on the subgroups becomes
$\prod_{uv\in E}c_{b_e}(y_u-y_v)$. A root coordinate together with the
oriented edge differences gives a bijective change of variables on a
tree. The differences are thus independent and uniform. Including the
densities of the subgroups, the magnitude of the original average is
\[
 p^{-\sum_v s_v}\prod_e d_e^{-1/2}=p^{-W(T,a)/2}
\]
by (2).

\paragraph{Freedom away from the subgroups.}
In fact \emph{every} unimodular filling of the zero entries gives the
same complex sum. To see this, take any proposed filling and multiply
each of its off-subgroup entries by an independent uniform sign. The
expected complex sum is the subgroup sum, by multilinearity. Each
realization lies in the closed disk of radius $p^{-W(T,a)/2}$ and its
mean lies on the boundary. All realizations must therefore coincide,
including the original filling. For odd $p$ the prescribed entries
are $p^k$-th roots, so filling every zero with one also proves the last
claim of Theorem~\ref{thm:ring}.

In the example modulo $27$, take $s_v=2$ at the center and zero at every
leaf. The construction prescribes the center phase only at $0,9,18$.
One can change its other $24$ values independently while retaining
exactly the same extremal sum.

\newpage
\section{A collision law over odd finite fields}
Let $q$ be an odd prime power, let every $a_{uv}\in\F_q^\times$, and put
\[
 F_g(x)=\sum_{uv\in E}a_{uv}x_ux_v+\sum_{v\in V}g_v(x_v),
 \qquad g_v:\F_q\longrightarrow\F_q.
\]
Write $P(T)$ for the minimum number of vertex-disjoint paths covering
$V$, allowing singleton paths, and $r=|V|-P(T)$. Equivalently, $r$
is the maximum number of edges retained in a linear forest in $T$.

\begin{theorem}\label{thm:collision}
For independent uniform $X,Y\in\F_q^V$,
\[
 \max_g\Pr[F_g(X)=F_g(Y)]
     =\frac1q+\left(1-\frac1q\right)q^{-r}. \tag{4}
\]
A single family of quadratic vertex functions attains the maximum
for every nontrivial additive character simultaneously.
\end{theorem}

\begin{proof}
For each nontrivial additive character $\chi$ of $\F_q$, the normalized
edge transform with kernel $\chi(a xy)$ has norm $q^{-1/2}$ whenever
$a\ne0$, by orthogonality. Path elimination gives
$|\E\chi(F_g(X))|\le q^{-r/2}$ for every $g$. Orthogonality again gives
\[
 \Pr[F_g(X)=F_g(Y)]=q^{-1}\sum_{\chi}\abs{\E\chi(F_g(X))}^2.
\]
For simultaneous attainment, apply (2) with every weight and $k$ equal
to one. It provides a subset $S$ with $2|S|+|E(T-S)|=r$ (take $S$ to
be the vertices with $s_v=1$). At $v\notin S$ choose
\[
  g_v(t)=-\tfrac12 t^2\sum_{u:\,uv\in E(T-S)}a_{uv},
  \qquad g_v=0\quad\text{on vertices in }S.
\]
Write the resulting form as $F_g(x)=\tfrac12 x^\mathsf{T}Bx$.
Its restriction to $x_S=0$ is
$-\tfrac12\sum_{uv\in E(T-S)}a_{uv}(x_u-x_v)^2$ and has rank
at most $|E(T-S)|$. Restoring $|S|$ rows and $|S|$ columns raises the
rank by at most $2|S|$, so $\operatorname{rank} B\le r$.
For each nontrivial $\chi$, differencing at $x+h$ and $x$ and averaging
in $x$ leaves exactly $h\in\ker B$. For such $h$ the residual phase is
one, so $|\E\chi(F_g(X))|^2=q^{-\operatorname{rank} B}$. The upper bound
forces $\operatorname{rank} B=r$ and gives equality in (4).
\end{proof}

For a star with $m\ge2$ arms of length two, one can retain all $m$
outer edges and exactly two inner edges. Thus $r=m+2$. In particular,
over $\F_3$ with four arms, the highest possible probability that two
outputs agree is $731/2187$, just $2/2187$ above the uniform baseline.
The maximum ranges over every choice of the nine local functions.

The tree hypothesis matters already on a triangle. Give its three
edges coefficient one and take $g_v(t)=t^2/2$. The resulting phase
polynomial is $(x_1+x_2+x_3)^2/2$, so its bias at a nontrivial character
has magnitude $q^{-1/2}$. This is sharp by conditioning on $x_3$ and
applying the single-edge estimate. The two edges of a spanning path
would have predicted $q^{-1}$ if the tree formula extended unchanged.

The equality between tree minimum rank and $|V|-P(T)$ is classical
\cite[Thm.~2.2]{chenette}. Its role here is visible in the quadratic
extremizer. The upper bound covers arbitrary vertex functions, and (1)
also accounts for interactions of different conductors. These are
research notes. Priority for the exact extremal formulas is unverified.

\begin{thebibliography}{9}
\bibitem{alpoge} L. Alp\"oge, \href{https://alpo.ge/remark.pdf}{announcement
on the Riemann hypothesis}, 20 November 2026.
\bibitem{schrijver} A. Schrijver,
\href{https://www-users.cse.umn.edu/~reiner/Classes/Schrijver2004.pdf}
{\emph{A Course in Combinatorial Optimization}}, sections on bipartite
matching and total unimodularity.
\bibitem{chenette} N. L. Chenette, S. V. Droms, L. Hogben, R. Mikkelson,
and O. Pryporova, \href{https://emis.de/ft/34661}{Minimum rank of a graph
over an arbitrary field}, \emph{Electronic Journal of Linear Algebra}
16 (2007), 183--186.
\end{thebibliography}
\end{document}
